\clearpage
\noindent\makebox[10mm][l]{42}\parbox[t]{\dimexpr\textwidth-10mm\relax}{\centering\scriptsize BESTIMMUNG DES BREITENUNTERSCHIEDS ZWISCHEN GÖTTINGEN UND ALTONA.}\par
1820 von der unteren Kulmination des 13. Mai an gerechnet:
\begin{center}Mai 13.\quad Obere Kulm.\ $-0''{,}10$;\end{center}
1824 von der oberen Kulmination des 20. April an gerechnet:
\[
\begin{array}{rl|r}
\text{Apr.}&21.\ \text{U. K.}&-0''{,}13\\
&21.\ \text{O. K.}&-0{,}26\\
&25.\ \text{U. K.}&-1{,}29\\
&27.\ \text{O. K.}&-2{,}04\\
&28.\ \text{O. K.}&-2{,}32\\
&29.\ \text{U. K.}&-2{,}45\\
\text{Mai}&1.\ \text{U. K.}&-2{,}93\\
&1.\ \text{O. K.}&-3{,}03\\
&2.\ \text{U. K.}&-3{,}14\\
&8.\ \text{O. K.}&-4{,}64\\
&9.\ \text{U. K.}&-4{,}77
\end{array}
\]
\begin{center}14.\end{center}
Bezeichnet man die Durchbiegung des Fernrohrs, das heißt die durch die Einwirkung der Schwerkraft auf sämtliche miteinander verbundenen Bestandteile des Instruments verursachte Lageänderung der auf die Ebene des geteilten Kreises projizierten optischen Achse gegenüber der Teilung, bei horizontaler optischer Achse mit $f$ und bei vertikaler mit $g$, und setzt voraus, dass diese Durchbiegung der Schwerkraft proportional ist (was angesichts der äußerst geringen Größe der gesamten Wirkung unbedenklich scheint), so wird die Durchbiegung bei einer Neigung $z$ der optischen Achse durch $f\sin z+g\cos z$ ausgedrückt. Dabei ist für den Kollimationsfehler $=e$ und die abgelesene Zenitdistanz $=z$ die wahre Zenitdistanz
\[
=z-e+f\sin(z-e)+g\cos(z-e)
\]
anzusetzen. Wäre das Fernrohr vollkommen symmetrisch, so würde $g$ ganz wegfallen. Da aber kein von Menschenhand gefertigter Gegenstand absolut vollkommen ist und überdies die Ausgleichsgewichte die vollkommene Symmetrie bereits in gewissem Maße stören, scheint es durchaus vernünftig, die Möglichkeit eines ein oder zwei Zehntel einer Sekunde betragenden

\clearpage
\noindent\mbox{}\hfill III.\quad ERGEBNISSE.\hfill 43\par
\noindent Wertes von $g$ zuzulassen. Wenn man die Rechnung auf einzelne Zehntel oder gar Hundertstel einer Sekunde genau führt, wäre es inkonsequent, den zweiten Teil der Durchbiegung unberücksichtigt zu lassen, soweit er berücksichtigt werden kann.
\begin{center}15.\end{center}
Die Ergänzung der halben Differenz zwischen der direkt und der durch Reflexion gemessenen Zenitdistanz zu $90^\circ$ ergibt die vom Kollimationsfehler und vom ersten Teil der Durchbiegung befreite Zenitdistanz. Sie enthält also nur noch den zweiten Teil der Durchbiegung, und zwar mit entgegengesetztem Vorzeichen, [je] nachdem der Kreis im Osten oder im Westen steht. Offenbar bezieht sich diese Zenitdistanz auf die Vertikale an der Stelle, an der die optische Achse das Wassergefäß trifft. Diese Vertikale ist für beide Kulminationen des Polarsterns unmerklich verschieden und liegt um $0''{,}05$ weiter nördlich als die Achse des Kreises. Diese Kombination ist für unseren Zweck auch insofern besser geeignet, als sie ohne die Voraussetzung auskommt, dass der Kollimationsfehler während der ganzen Dauer der Beobachtungen von 1824 unverändert blieb. Setzt man die Zahl der direkten Beobachtungen $=\alpha$ und die der Reflexionsbeobachtungen $=\beta$, so ist das Gewicht jener Bestimmung $=\dfrac{4\alpha\beta}{\alpha+\beta}$, sofern man die Beobachtungsfehler als rein zufällig und voneinander unabhängig ansieht.
\begin{center}16.\end{center}
Bezeichnen wir nun mit
\begin{quote}
\begin{tabular}{@{}l@{\quad}p{.84\textwidth}@{}}
$\varphi$&die Polhöhe am Standort des Wassergefäßes\\
$\delta$&die Deklination des Polarsterns bei der unteren Kulmination am 13. Mai 1820\\
$\delta'$&die Deklination bei der oberen Kulmination am 20. April 1824,
\end{tabular}
\end{quote}
so geben uns die Beobachtungen folgende Bestimmungen:
\begin{center}
\begin{tabular}{lllcll}
\multicolumn{6}{l}{für $\delta+\varphi-0{,}765g$}\\
\quad1820&Mai 13&$\ldots\,\ldotp$&$139^\circ52'38''{,}40$&Gewicht&$6{,}86$\\
\multicolumn{6}{l}{für $\delta-\varphi+0{,}800g$}\\
\quad1820&Mai 13&$\ldots\,\ldotp$&$36^\circ49'\phantom0 1''{,}50$&Gewicht&$2{,}00$
\end{tabular}
\end{center}
\noindent\hfill $6^*$\par

\clearpage
\noindent\makebox[10mm][l]{44}\parbox[t]{\dimexpr\textwidth-10mm\relax}{\centering\scriptsize BESTIMMUNG DES BREITENUNTERSCHIEDS ZWISCHEN GÖTTINGEN UND ALTONA.}\par
\begin{center}
\begin{tabular}{rrrc rr r c r}
\multicolumn{9}{l}{für $\delta'+\varphi-0{,}765g$}\\
1824&Apr.&21&$\ldots\,\ldotp$&$139^\circ$&$54'$&$5''{,}61$&Gewicht&$6{,}86$\\
&&25&$\ldots\,\ldotp$&139&54&$5{,}76$&»&$6{,}86$\\
&&29&$\ldots\,\ldotp$&139&54&$6{,}36$&»&$6{,}86$\\
&Mai&1&$\ldots\,\ldotp$&139&54&$5{,}91$&»&$6{,}86$\\
\multicolumn{9}{l}{für $\delta'-\varphi+0{,}800g$}\\
1824&Apr.&20&$\ldots\,\ldotp$&$36^\circ$&$50'$&$31''{,}15$&Gewicht&$2{,}67$\\
&&21&$\ldots\,\ldotp$&36&50&$30{,}29$&»&$6{,}86$\\
&&27&$\ldots\,\ldotp$&36&50&$30{,}65$&»&$6{,}86$\\
&&28&$\ldots\,\ldotp$&36&50&$30{,}73$&»&$6{,}86$\\
&Mai&1&$\ldots\,\ldotp$&36&50&$30{,}25$&»&$6{,}86$\\
\multicolumn{9}{l}{für $\delta'+\varphi+0{,}765g$}\\
1824&Mai&2&$\ldots\,\ldotp$&$139^\circ$&$54'$&$6''{,}71$&Gewicht&$6{,}86$\\
&&9&$\ldots\,\ldotp$&139&54&$6{,}15$&»&$6{,}86$\\
\multicolumn{9}{l}{für $\delta'-\varphi-0{,}800g$}\\
1824&Mai&8&$\ldots\,\ldotp$&$36^\circ$&$50'$&$30''{,}48$&Gewicht&$6{,}86$.
\end{tabular}
\end{center}
Wir erhalten demnach zur Bestimmung der vier unbekannten Größen $\delta$, $\delta'$, $\varphi$, $g$ die sechs Gleichungen
\[
\begin{array}{rclrrrcr}
\delta+\varphi-0{,}765g&=&139^\circ&52'&38''{,}40&&\text{Gewicht}&6{,}86\\
\delta-\varphi+0{,}800g&=&36&49&1{,}50&&\text{»}&2{,}00\\
\delta'+\varphi-0{,}765g&=&139&54&5{,}91&&\text{»}&27{,}43\\
\delta'-\varphi+0{,}800g&=&36&50&30{,}54&&\text{»}&30{,}10\\
\delta'+\varphi+0{,}765g&=&139&54&6{,}43&&\text{»}&13{,}71\\
\delta'-\varphi-0{,}800g&=&36&50&30{,}48&&\text{»}&6{,}86.
\end{array}
\]
woraus sich nach der Methode der kleinsten Quadrate\textsuperscript{*)} folgende Werte ergeben:
\[
\begin{array}{rcrrr}
\delta&=&88^\circ&20'&50''{,}33\\
\delta'&=&88&22&18{,}28\\
\varphi&=&51&31&47{,}90\\
g&=&\multicolumn{3}{r}{+\ 0{,}17.}
\end{array}
\]
Das Gewicht der Bestimmung von $\varphi$ beträgt dabei $=60{,}8$.
\par\noindent\rule{.25\textwidth}{.4pt}\par
{\footnotesize *) Hier etwas bequemer nach dem Verfahren in Supplem.\ Theor.\ Comb.\ Observ.\par}

\clearpage
\noindent\mbox{}\hfill III.\quad ERGEBNISSE.\hfill 45\par
Um einen ungefähren Maßstab für die Genauigkeit der Beobachtungen zu gewinnen, setzen wir diese Werte in die vierzehn Gleichungen ein, aus denen die vorherigen sechs zusammengefasst worden waren. Es bleiben dann folgende Fehler übrig:
\begin{center}
\begin{tabular}{r|c}
Fehler&Gewicht der Gleichung\\\hline
$-0''{,}30$&$6{,}86$\\
$+1{,}07$&$2{,}00$\\
$+0{,}44$&$6{,}86$\\
$+0{,}29$&$6{,}86$\\
$-0{,}31$&$6{,}86$\\
$+0{,}14$&$6{,}86$\\
$-0{,}63$&$2{,}67$\\
$+0{,}23$&$6{,}86$\\
$-0{,}13$&$6{,}86$\\
$-0{,}21$&$6{,}86$\\
$+0{,}27$&$6{,}86$\\
$-0{,}40$&$6{,}86$\\
$+0{,}16$&$6{,}86$\\
$-0{,}24$&$6{,}86$
\end{tabular}
\end{center}
Die Summe der Produkte aus den Quadraten dieser Fehler und den Gewichten ist $=9{,}6154$. Somit ergibt sich als Näherungswert für den mittleren Fehler einer Beobachtung
\[
=\sqrt{\frac{9{,}6154}{10}}=0''{,}981.
\]
Der im Endergebnis für die Polhöhe zu befürchtende mittlere Fehler ist demnach, soweit er von unregelmäßig wirkenden Ursachen herrührt,
\[
=\frac{0''{,}981}{\sqrt{60{,}8}}=0''{,}126.
\]
Die Unsicherheit des Ergebnisses muss allerdings etwas größer sein, da die Annahme, dass sämtliche Beobachtungsfehler ohne erkennbare Ordnung und voneinander unabhängig auftreten, nicht ganz richtig ist. Bei gleichartigen Beobach-

\clearpage
\noindent\makebox[10mm][l]{46}\parbox[t]{\dimexpr\textwidth-10mm\relax}{\centering\scriptsize BESTIMMUNG DES BREITENUNTERSCHIEDS ZWISCHEN GÖTTINGEN UND ALTONA.}\par
\noindent tungen zu einer Kulmination und bei gleichartigen Kulminationen an mehreren Tagen liegt nämlich annähernd dasselbe Ableseergebnis zugrunde. Bei der Ablesung mit Nonien kommen zwar fast immer andere Teilstriche zum Tragen, sodass deren unregelmäßige Teilungsfehler bei unserem Verfahren im mittleren Fehler einer Beobachtung $0''{,}981$ enthalten sind. Dennoch ist es natürlich, dass in den verschiedenen Bereichen des Teilkreises unterschiedliche Durchschnittsfehler vorherrschen müssen. Sie sind jedenfalls sehr klein. Im Jahr 1826 habe ich mit vier ausgezeichneten Mikroskopen von \textsc{Repsold} 30 Teilstriche von 12 zu 12 Grad mit äußerster Sorgfalt geprüft. Dabei wurde jeder Teilstrich in wechselnden Kombinationen fast 200-mal eingestellt. Das Ergebnis ist, dass der Mittelwert der Fehler zweier diametral gegenüberliegender Teilstriche, $A$ und $A+180^\circ$, soweit sich noch eine gewisse Regelmäßigkeit erkennen lässt, möglichst genau durch die Formel
\[
-1''{,}23\cos(2A-28^\circ28')-0''{,}22\cos(4A-47^\circ56')
\]
dargestellt wird. Die dann verbleibenden Fehler erscheinen regellos, und die Quadratwurzel aus dem Mittelwert ihrer Quadrate ist $=0''{,}32$. Ich hatte mir vorgenommen, diese Prüfung auf die doppelte Anzahl von Teilstrichen auszudehnen. Angesichts der Geringfügigkeit der Ergebnisse scheint die Untersuchung jedoch den großen Zeitaufwand, den sie erfordern würde, nicht zu rechtfertigen. Der erste Teil des regelmäßigen Fehlers $-1''{,}23\cos(2A-28^\circ28')$ fällt selbstverständlich von selbst weg, wenn man, wie bei den obigen Beobachtungen stets geschehen, alle vier Nonien abliest. Er liefert hingegen eine tatsächliche Korrektur, falls man die Teilung nur an zwei gegenüberliegenden Stellen abliest, wie ich es jetzt immer tue, seit ich anstelle der Nonien zwei \textsc{Repsold}sche Mikroskope verwende und dadurch erheblich an Ablesegenauigkeit gewonnen habe.
\begin{center}17.\end{center}
Zieht man es vor, $g=0$ vorauszusetzen, so fällt die Polhöhe um $0''{,}07$ kleiner aus, und das Gewicht dieser Bestimmung beträgt $=84{,}1$. Weitere, an anderer Stelle anzuführende Beobachtungen scheinen übrigens den obigen

\clearpage
\noindent\mbox{}\hfill III.\quad ERGEBNISSE.\hfill 47\par
\noindent Wert von $g$ im Vorzeichen und sehr annähernd auch in der Größe zu bestätigen, reichen aber noch nicht aus, um eine so heikle Frage zu entscheiden.

Den Koeffizienten $f$ kann man aus den vorliegenden Beobachtungen nur unter der Voraussetzung bestimmen, dass der Kollimationsfehler während der Beobachtungen von 1824 unverändert blieb. Lässt man diese Voraussetzung zu, so hat man 28 Gleichungen, deren sachgerechte Behandlung
\[
\begin{array}{rcll}
\varphi&=&51^\circ31'47''{,}89&\text{mit dem Gewicht }60{,}9\\
f&=&\phantom{51^\circ31'}+\ 0{,}76&\\
g&=&\phantom{51^\circ31'}+\ 0{,}23&
\end{array}
\]
ergibt. Da man den Kollimationsfehler heute durch Einstellen des Fernrohrs auf den Nadirpunkt zu jeder Stunde mit erstaunlicher Genauigkeit ohne Umlegen bestimmen kann\textsuperscript{*)}, behalte ich mir eine weitere Prüfung dieser Frage vor.
\begin{center}18.\end{center}
Unter dem Vorbehalt einer durch weitere künftige Untersuchungen noch zu ermittelnden Korrektur, die wohl kaum eine halbe Sekunde erreichen kann, setze ich daher die Polhöhe an
\begin{center}
\begin{tabular}{p{.72\textwidth}r}
in Göttingen&\\
\quad am Standort des Wassergefäßes bei den Beobachtungen des Polarsterns&$51^\circ31'47''{,}90$\\
\quad am Standort des \textsc{Reichenbach}schen Meridiankreises&$47{,}85$\\
\quad am Standort des Zenitsektors&$47{,}92$\\
\multicolumn{2}{p{.94\textwidth}}{\quad\quad (Letztere wurde der Reduktion der Deklinationen der Zenitsterne zugrunde gelegt)}\\
in Altona&\\
\quad am Standort des Zenitsektors&$53^\circ32'44''{,}42$\\
\quad am Standort des Meridiankreises&$45{,}27$.
\end{tabular}
\end{center}
\par\noindent\rule{.25\textwidth}{.4pt}\par
{\footnotesize *) Ich benutze dieses unschätzbare Mittel, dessen Durchführbarkeit \textsc{Bohnenberger} als Erster gezeigt hat, seit zwei Jahren ständig.\par}

\clearpage
\noindent\makebox[10mm][l]{48}\parbox[t]{\dimexpr\textwidth-10mm\relax}{\centering\scriptsize BESTIMMUNG DES BREITENUNTERSCHIEDS ZWISCHEN GÖTTINGEN UND ALTONA.}\par
\begin{center}19.\end{center}
Nach der trigonometrischen Verbindung der Sternwarten von Göttingen und Altona liegt letztere
\[
\begin{array}{rl}
115163{,}725&\text{Toisen nördlich}\\
7{,}211&\text{Toisen westlich}
\end{array}
\]
von der ersteren. Diese Zahlen beziehen sich auf die Standorte der Meridiankreise. Sie beruhen auf dem Wert der Dreiecksseite Hamburg-Hohenhorn von $13841{,}815$ Toisen, der wiederum auf der von Herrn\ Prof.\ \textsc{Schumacher} im Jahr 1820 in Holstein gemessenen Basis beruht. Da aber der Vergleich der dabei verwendeten Messstangen mit der Normaltoise noch nicht \textls[150]{endgültig} abgeschlossen ist, wird sich die obige Entfernung künftig noch im selben Verhältnis ändern wie die Basis selbst. Diese Änderung kann jedenfalls nur sehr gering sein. Die mittlere Länge eines Breitengrades zwischen beiden Sternwarten ergibt sich demnach zu
\[
=57127{,}2\text{ Toisen},
\]
merklich größer, als man nach den mittleren Werten der in Frankreich und England gemessenen Grade hätte erwarten können.
\begin{center}20.\end{center}
Die hannoversche Gradmessung liefert also einen weiteren Beitrag zur Bestätigung der nicht mehr zu bezweifelnden Tatsache, dass die Erdoberfläche keine ganz regelmäßige Gestalt hat. Die Anomalien in den Teilabschnitten der französischen und der englischen Gradmessung haben diese Unregelmäßigkeit bereits belegt. Noch stärkere Belege dafür haben die Anomalien bei den Polhöhen mehrerer Orte in Italien geliefert. Bei der hannoverschen Gradmessung zeigt sich neben der Anomalie zwischen Göttingen und Altona eine noch beträchtlich stärkere an einem dazwischenliegenden Dreieckspunkt, dem Brocken. Berechnet man meine Dreiecke so, als lägen sie auf der Oberfläche eines elliptischen Sphäroids, dessen Abmessungen \textsc{Walbeck} aus der Gesamtheit der bisherigen Gradmessungen abgeleitet hat und das nach unserem besten heutigen Wissen der wirklichen Gestalt

\clearpage
\noindent\mbox{}\hfill III.\quad ERGEBNISSE.\hfill 49\par
\noindent\textls[150]{im Ganzen} am nächsten kommt (Abplattung $\dfrac1{302{,}78}$, der dreihundertsechzigste Teil des Erdmeridians $=57009{,}758$ Toisen), und geht man dabei von der Polhöhe Göttingens $=51^\circ31'47''{,}85$ aus, so erhält man die Breite
\[
\begin{array}{lcrrr}
\text{des Brockens}&=&51^\circ&48'&1''{,}85\\
\text{von Altona}&=&53&32&50{,}79.
\end{array}
\]
Während nun die astronomischen Beobachtungen für Altona eine um $5''{,}52$ \textls[150]{kleinere} Polhöhe ergeben haben, ergibt sich aus den von Herrn\ von \textsc{Zach} auf dem Brocken angestellten Beobachtungen für diesen Punkt eine um $10\text{—}11''$ \textls[150]{größere}\textsuperscript{*)} Polhöhe. Von diesem Unterschied kann jedenfalls nur ein kleiner Teil dem Instrument und den bei der Rechnung verwendeten Deklinationen zugeschrieben werden. Vergliche man den Breitenunterschied zwischen Altona und dem Brocken mit der Krümmung, die dem am besten an die Erde im Ganzen angepassten Sphäroid entspricht, so ergäbe sich daher eine Abweichung von $16''$.

Unseres Erachtens betrachtet man diese Frage aus dem falschen Blickwinkel, wenn man bei solchen Erscheinungen stets nur von lokalen Ablenkungen der Lotlinie spricht und sie damit gleichsam als einzelne Ausnahmen ansieht. Was wir im geometrischen Sinn die Erdoberfläche nennen, ist genau die Fläche, welche die Richtung der Schwerkraft überall senkrecht schneidet und von der die Oberfläche des Weltmeers einen Teil bildet. Die Richtung der Schwerkraft an jedem Punkt wird durch die Gestalt des festen Teils der Erde und seine ungleichmäßige Dichte bestimmt. An der äußeren Kruste der Erde, über die allein wir etwas wissen, erweisen sich diese Gestalt und Dichte als höchst unregelmäßig. Die Unregelmäßigkeit der Dichte kann sich durchaus noch ziemlich weit unter die äußere Kruste erstrecken und entzieht sich ganz unseren Berechnungen, für die fast alle Daten fehlen. Die geometrische Oberfläche ist das Produkt der Gesamtwirkung dieser ungleich verteilten Elemente. Wenn sich eindeutige Belege für die Unregelmäßigkeit finden, sollte man sich weniger darüber wundern, als vielmehr darüber, dass sie nicht noch größer ist.
\par\noindent\rule{.25\textwidth}{.4pt}\par
{\footnotesize *) Monatl.\ Corresp.\ B.\ X.\ S.\ 203. An einer Stelle, die etwa $0''{,}5$ südlicher liegt als der Dreieckspunkt, ermittelte dieser geschickte Beobachter aus 188 Beobachtungen von $\alpha$ Aquilae $51^\circ48'12''{,}12$. Aus Sonnenbeobachtungen ermittelte er $51^\circ48'11''{,}17$.\par}
\noindent IX.\hfill 7\par

\clearpage
\noindent\makebox[10mm][l]{50}\parbox[t]{\dimexpr\textwidth-10mm\relax}{\centering\scriptsize BESTIMMUNG DES BREITENUNTERSCHIEDS ZWISCHEN GÖTTINGEN UND ALTONA.}\par
\noindent Ließe sich bei den astronomischen Beobachtungen eine zehn- oder hundertmal größere Genauigkeit als heute erreichen, so würden sie diese Unregelmäßigkeit zweifellos überall nachweisen.

Bei dieser Sachlage hindert uns jedoch nichts, die Erde im Ganzen als ein elliptisches Rotationssphäroid zu betrachten, von dem die wirkliche (geometrische) Oberfläche überall in bald stärkeren, bald schwächeren, bald kürzeren, bald längeren Wellen abweicht. Könnte man die ganze Erde gleichsam mit einem einzigen trigonometrischen Netz umspannen und daraus die gegenseitige Lage aller Punkte berechnen, so wäre das ideale Rotationssphäroid jenes, auf dem die berechneten Richtungen der Vertikalen möglichst gut mit den astronomischen Beobachtungen übereinstimmten. Auch wenn wir von diesem unerreichbaren Ideal stets weit entfernt bleiben werden, besteht doch kein Zweifel, dass die kommenden Jahrhunderte die mathematische Kenntnis der Erdgestalt erheblich voranbringen können. Die Vermehrung der Gradmessungen ist aber eigentlich nur der Anfang. Sie liefert nur einzelne Ergebnisse für eine kleine Anzahl von Punkten auf voneinander getrennten Linien. Wie viel ergiebiger wird die Ausbeute sein, wenn die in verschiedenen Ländern mit sorgfältig ausgewählten Hilfsmitteln ausgeführten trigonometrischen Arbeiten miteinander verknüpft und zu einem einzigen großen System zusammengeführt werden. Vielleicht ist die Aussicht, dass eines Tages alle Sternwarten Europas trigonometrisch miteinander verbunden sein werden, keine bloße Wunschvorstellung. Schon jetzt bestehen solche Verbindungen von Schottland bis zum Adriatischen Meer und von Formentera bis Fünen, auch wenn sie bisher nur teilweise öffentlich bekannt gemacht worden sind. Möge doch dieser letzte Umstand mehr Beachtung finden als bisher, und mögen die kostbaren Materialien, die der wissenschaftlichen Welt gehören sollten, ihr weder vorenthalten noch gar der Gefahr ausgesetzt werden, verloren zu gehen!
\begin{center}21.\end{center}
Ein nicht uninteressantes Ergebnis liefert auch der Vergleich der aus den Sektorbeobachtungen gewonnenen Sterndeklinationen mit älteren Bestimmungen, soweit solche vorliegen. Von unseren 43 Sternen finden sich

\clearpage
\noindent\mbox{}\hfill III.\quad ERGEBNISSE.\hfill 51\par
\noindent 27 im Katalog von \textsc{Piazzi} und 13 in \textsc{Bessel}s \textsc{Bradley}-Katalog. Es folgt der Vergleich unserer Bestimmungen (1827) mit \raisebox{.3ex}{.}denen von \textsc{Bradley} (1755) und \textsc{Piazzi} (1800), reduziert nach \textsc{Bessel}s neuer Bestimmung der Präzession. Positive Vorzeichen bedeuten nach unserer Bestimmung eine weiter nördliche Stellung.
\begin{center}
{\small\setlength{\tabcolsep}{4pt}
\begin{tabular}{r|l|r|r||r|l|r|r}
\multicolumn{2}{c|}{Bezeichnung}&\textsc{Bradley}&\textsc{Piazzi}&\multicolumn{2}{c|}{Bezeichnung}&\textsc{Bradley}&\textsc{Piazzi}\\\cline{1-4}\cline{5-8}
1&24 Canum&$+0''{,}2$&$+1''{,}1$&17&P.\ 15.39&—&$-0''{,}7$\\
2&83 Ursae&$-1{,}5$&$-2{,}0$&25&$\theta$ Draconis&$+23''{,}5$&$+8{,}6$\\
3&$\eta$ Ursae&$-2{,}4$&$-2{,}3$&27&P.\ 16.33&—&$-2{,}6$\\
4&86 Ursae&$-5{,}4$&$-0{,}8$&28&P.\ 16.56&—&$-2{,}6$\\
6&P.\ 13.289&—&$-1{,}3$&32&16 Draconis&$+1{,}2$&$-3{,}7$\\
7&13 Bootis&$+1{,}4$&$+2{,}4$&36&P.\ 16.253&—&$-2{,}1$\\
8&$\varkappa$ Bootis sq.&$-2{,}9$&$-2{,}2$&37&P.\ 16.291&—&$+10{,}3$\\
9&P.\ 14.56&—&$-0{,}4$&38&P.\ 16.310&—&$-4{,}9$\\
10&$\theta$ Bootis&$-30{,}6$&$-10{,}8$&39&P.\ 17.20&—&$-3{,}0$\\
11&P.\ 14.131&—&$+7{,}4$&40&P.\ 17.38&—&$-2{,}9$\\
12&P.\ 14.164&—&$+0{,}1$&41&74 Herculis&$+1{,}7$&$+0{,}7$\\
13&39 Bootis med.&$+4{,}8$&$+2{,}7$&42&P.\ 17.120&—&$-2{,}1$\\
14&P.\ 14.235&—&$-5{,}0$&43&$\beta$ Draconis&$-0{,}5$&$-1{,}5$\\
15&44 Bootis med.&$+1{,}9$&$+1{,}1$&&&&
\end{tabular}}
\end{center}
\vspace{14mm}
\begin{center}\rule{.25\textwidth}{.4pt}\end{center}
\noindent\hfill $7^*$\par
